\subsection{\label{sec.details.domino}Domino tilings}

\begin{fineprint}
In Subsection \ref{subsec.gf.weighted-set.domino}, we have given a
classification of faultfree domino tilings of the rectangle $R_{n,3}$. Let us
now outline how this classification can be proved. First, we introduce names
for the faultfree domino tilings that appear in our classification:

\begin{definition}
\label{def.gf.weighted-set.domino.Rn3.ABC}\textbf{(a)} For each even positive
integer $n$, we let $A_{n}$ be the domino tiling%
\[%
%TCIMACRO{\TeXButton{tikz tilings An}{\qquad% [tikz figure removed]}}%
%BeginExpansion
\qquad% [tikz figure removed]%
%EndExpansion
\ \ \ \ \ \ \ \ \ \ \text{of }R_{n,3}.
\]
Formally speaking, this domino tiling is the set partition of $R_{n,3}$
consisting of the following the dominos:

\begin{itemize}
\item the horizontal dominos $\left\{  \left(  2i-1,\ 1\right)  ,\ \left(
2i,\ 1\right)  \right\}  $ for all $i\in\left[  n/2\right]  $, which fill the
bottom row of $R_{n,3}$, and which we call the \textbf{basement dominos};

\item the vertical domino $\left\{  \left(  1,2\right)  ,\ \left(  1,3\right)
\right\}  $ in the first column, which we call the \textbf{left wall};

\item the vertical domino $\left\{  \left(  n,2\right)  ,\ \left(  n,3\right)
\right\}  $ in the last column, which we call the \textbf{right wall};

\item the horizontal dominos $\left\{  \left(  2i,\ 2\right)  ,\ \left(
2i+1,\ 2\right)  \right\}  $ for all $i\in\left[  n/2-1\right]  $, which fill
the middle row of $R_{n,3}$ (except for the first and last columns), and which
we call the \textbf{middle dominos};

\item the horizontal dominos $\left\{  \left(  2i,\ 3\right)  ,\ \left(
2i+1,\ 3\right)  \right\}  $ for all $i\in\left[  n/2-1\right]  $, which fill
the top row of $R_{n,3}$ (except for the first and last columns), and which we
call the \textbf{top dominos}.
\end{itemize}

\textbf{(b)} For each even positive integer $n$, we let $B_{n}$ be the domino
tiling%
\[%
%TCIMACRO{\TeXButton{tikz tilings Bn}{\qquad% [tikz figure removed]}}%
%BeginExpansion
\qquad% [tikz figure removed]%
%EndExpansion
\ \ \ \ \ \ \ \ \ \ \text{of }R_{n,3}.
\]
This is the reflection of the tiling $A_{n}$ across the horizontal axis of
symmetry of $R_{n,3}$. \medskip

\textbf{(c)} We let $C$ denote the domino tiling
\[%
% [tikz figure removed]
%BeginExpansion
% [tikz figure removed]%
%EndExpansion
\ \ \ \ \ \ \ \ \ \ \text{of }R_{2,3}.
\]

\end{definition}

Our classification now can be stated as follows:

\begin{proposition}
\label{prop.gf.weighted-set.domino.Rn3.ABC}The faultfree domino tilings of
height-$3$ rectangles are precisely the tilings%
\[
A_{2},A_{4},A_{6},A_{8},\ldots,\ \ \ \ \ \ \ \ \ \ B_{2},B_{4},B_{6}%
,B_{8},\ldots,\ \ \ \ \ \ \ \ \ \ C
\]
we have just defined. More concretely: \medskip

\textbf{(a)} The faultfree domino tilings of a height-$3$ rectangle that
contain a vertical domino in the \textbf{top} two squares of the first column
are $A_{2},A_{4},A_{6},A_{8},\ldots$. \medskip

\textbf{(b)} The faultfree domino tilings of a height-$3$ rectangle that
contain a vertical domino in the \textbf{bottom} two squares of the first
column are $B_{2},B_{4},B_{6},B_{8},\ldots$. \medskip

\textbf{(c)} The only faultfree domino tiling of a height-$3$ rectangle that
contains \textbf{no} vertical domino in the first column is $C$.
\end{proposition}

\begin{proof}
[Proof of Proposition \ref{prop.gf.weighted-set.domino.Rn3.ABC} (sketched).]It
clearly suffices to prove parts \textbf{(a)}, \textbf{(b)} and \textbf{(c)}.
We begin with the easiest part, which is \textbf{(c)}: \medskip

\textbf{(c)} Clearly, $C$ is a faultfree domino tiling of $R_{2,3}$ that
contains no vertical domino in the first column. It remains to show that $C$
is the only such tiling.

Indeed, let $T$ be any faultfree domino tiling of $R_{2,3}$ that contains no
vertical domino in the first column. Thus, its first column must be filled
with three horizontal dominos, which all must protrude into the second column
and thus cover that second column as well. If $T$ had any further column, then
$T$ would have a fault between its $2$-nd and its $3$-rd column, which is
impossible for a faultfree tiling. Thus, $T$ must consist only of the three
horizontal dominos already mentioned. In other words, $T=C$. This completes
the proof of Proposition \ref{prop.gf.weighted-set.domino.Rn3.ABC}
\textbf{(c)}. \medskip

\textbf{(a)} It is straightforward to check that the tilings $A_{2}%
,A_{4},A_{6},A_{8},\ldots$ are faultfree (indeed, the basement dominos prevent
faults between the $\left(  2i-1\right)  $-st and $\left(  2i\right)  $-th
columns, whereas the top dominos prevent faults between the $\left(
2i\right)  $-th and $\left(  2i+1\right)  $-st columns). Thus, they are
faultfree domino tilings of a height-$3$ rectangle that contain a vertical
domino in the \textbf{top} two squares of the first column. It remains to show
that they are the only such tilings.

Indeed, let $T$ be any faultfree domino tiling of a height-$3$ rectangle that
contains a vertical domino in the \textbf{top} two squares of the first
column. Let $n$ be the width of this rectangle (so that the rectangle is
$R_{n,3}$). We shall show that $n$ is even, and that $T=A_{n}$.

We know that $T$ contains a vertical domino in the \textbf{top} two squares of
the first column. In other words, $T$ contains the left wall (where we are
using the terminology from Definition \ref{def.gf.weighted-set.domino.Rn3.ABC}
\textbf{(a)}). The remaining square in the first column is the square $\left(
1,1\right)  $, and it must thus be covered by the basement domino $\left\{
\left(  1,1\right)  ,\ \left(  2,1\right)  \right\}  $ (since no other domino
would fit). Hence, $T$ contains the basement domino $\left\{  \left(
1,1\right)  ,\ \left(  2,1\right)  \right\}  $.

We shall now prove the following claims:

\begin{statement}
\textit{Claim 1:} For each positive integer $i<n/2$, the tiling $T$ contains
the basement domino $\left\{  \left(  2i-1,\ 1\right)  ,\ \left(
2i,\ 1\right)  \right\}  $, the middle domino $\left\{  \left(  2i,\ 2\right)
,\ \left(  2i+1,\ 2\right)  \right\}  $ and the top domino $\left\{  \left(
2i,\ 3\right)  ,\ \left(  2i+1,\ 3\right)  \right\}  $.
\end{statement}

\begin{proof}
[Proof of Claim 1.]We induct on $i$:

\textit{Base case:} We must prove that Claim 1 holds for $i=1$, provided that
$1<n/2$. So let us assume that $1<n/2$. Thus, $2<n$, so that $n\geq3$. We must
prove that Claim 1 holds for $i=1$; in other words, we must prove that $T$
contains the basement domino $\left\{  \left(  1,1\right)  ,\ \left(
2,1\right)  \right\}  $, the middle domino $\left\{  \left(  2,2\right)
,\ \left(  3,2\right)  \right\}  $ and the top domino $\left\{  \left(
2,3\right)  ,\ \left(  3,3\right)  \right\}  $. For the basement domino, we
have already proved this. For the middle and the top domino, we argue as
follows: If $T$ had a vertical domino in the $2$-nd column, then this domino
would cover the top two squares of that column (since the bottom square is
already covered by the basement domino $\left\{  \left(  1,1\right)
,\ \left(  2,1\right)  \right\}  $), and thus $T$ would have a fault between
the $2$-nd and $3$-rd columns (since the basement domino $\left\{  \left(
1,1\right)  ,\ \left(  2,1\right)  \right\}  $ also ends at the $2$-nd
column), which would contradict the faultfreeness of $T$. Hence, $T$ has no
vertical domino in the $2$-nd column. Thus, the top two squares of the $2$-nd
column of $T$ must be covered by horizontal dominos. These horizontal dominos
must both protrude into the $3$-rd column (since the corresponding squares in
the $1$-st column are already covered by the left wall), and thus must be the
middle domino $\left\{  \left(  2,2\right)  ,\ \left(  3,2\right)  \right\}  $
and the top domino $\left\{  \left(  2,3\right)  ,\ \left(  3,3\right)
\right\}  $. Hence, we have shown that $T$ contains the middle domino
$\left\{  \left(  2,2\right)  ,\ \left(  3,2\right)  \right\}  $ and the top
domino $\left\{  \left(  2,3\right)  ,\ \left(  3,3\right)  \right\}  $. This
completes the base case.

\textit{Induction step:} Let $j$ be a positive integer such that $j+1<n/2$.
Assume (as the induction hypothesis) that Claim 1 holds for $i=j$. In other
words, $T$ contains the basement domino $\left\{  \left(  2j-1,\ 1\right)
,\ \left(  2j,\ 1\right)  \right\}  $, the middle domino $\left\{  \left(
2j,\ 2\right)  ,\ \left(  2j+1,\ 2\right)  \right\}  $ and the top domino
$\left\{  \left(  2j,\ 3\right)  ,\ \left(  2j+1,\ 3\right)  \right\}  $. We
must now show that Claim 1 holds for $i=j+1$ as well, i.e., that $T$ also
contains the basement domino $\left\{  \left(  2j+1,\ 1\right)  ,\ \left(
2j+2,\ 1\right)  \right\}  $, the middle domino $\left\{  \left(
2j+2,\ 2\right)  ,\ \left(  2j+3,\ 2\right)  \right\}  $ and the top domino
$\left\{  \left(  2j+2,\ 3\right)  ,\ \left(  2j+3,\ 3\right)  \right\}  $.

Indeed, we first recall that $j+1<n/2$, so that $2\left(  j+1\right)  <n$ and
therefore $n>2\left(  j+1\right)  =2j+2$, so that $n\geq2j+3$. This shows that
the rectangle $R_{n,3}$ has a $\left(  2j+3\right)  $-th column (along with
all columns to its left).

Now, the square $\left(  2j+1,\ 1\right)  $ cannot be covered by a vertical
domino in $T$, since this vertical domino would collide with the middle domino
$\left\{  \left(  2j,\ 2\right)  ,\ \left(  2j+1,\ 2\right)  \right\}  $
(which we already know to belong to $T$). Thus, this square must be covered by
a horizontal domino. This horizontal domino cannot protrude into the $\left(
2j\right)  $-th column (since it would then collide with the basement domino
$\left\{  \left(  2j-1,\ 1\right)  ,\ \left(  2j,\ 1\right)  \right\}  $,
which we already know to belong to $T$), and thus must be the basement domino
$\left\{  \left(  2j+1,\ 1\right)  ,\ \left(  2j+2,\ 1\right)  \right\}  $. So
we have shown that $T$ contains the basement domino $\left\{  \left(
2j+1,\ 1\right)  ,\ \left(  2j+2,\ 1\right)  \right\}  $. It remains to show
that $T$ also contains the middle domino $\left\{  \left(  2j+2,\ 2\right)
,\ \left(  2j+3,\ 2\right)  \right\}  $ and the top domino $\left\{  \left(
2j+2,\ 3\right)  ,\ \left(  2j+3,\ 3\right)  \right\}  $.

If $T$ had a vertical domino in the $\left(  2j+2\right)  $-nd column, then
this domino would cover the top two squares of that column (since the bottom
square is already covered by the basement domino $\left\{  \left(
2j+1,\ 1\right)  ,\ \left(  2j+2,\ 1\right)  \right\}  $), and thus $T$ would
have a fault between the $\left(  2j+2\right)  $-nd and $\left(  2j+3\right)
$-rd columns (since the basement domino $\left\{  \left(  2j+1,\ 1\right)
,\ \left(  2j+2,\ 1\right)  \right\}  $ also ends at the $\left(  2j+2\right)
$-nd column), which would contradict the faultfreeness of $T$. Hence, $T$ has
no vertical domino in the $\left(  2j+2\right)  $-nd column. Thus, the top two
squares of the $\left(  2j+2\right)  $-nd column of $T$ must be covered by
horizontal dominos. These horizontal dominos must both protrude into the
$\left(  2j+3\right)  $-rd column (since the corresponding squares in the
$\left(  2j+1\right)  $-st column are already covered by the middle domino
$\left\{  \left(  2j,\ 2\right)  ,\ \left(  2j+1,\ 2\right)  \right\}  $ and
the top domino $\left\{  \left(  2j,\ 3\right)  ,\ \left(  2j+1,\ 3\right)
\right\}  $), and thus must be the middle domino $\left\{  \left(
2j+2,\ 2\right)  ,\ \left(  2j+3,\ 2\right)  \right\}  $ and the top domino
$\left\{  \left(  2j+2,\ 3\right)  ,\ \left(  2j+3,\ 3\right)  \right\}  $.
Hence, we have shown that $T$ contains the middle domino $\left\{  \left(
2j+2,\ 2\right)  ,\ \left(  2j+3,\ 2\right)  \right\}  $ and the top domino
$\left\{  \left(  2j+2,\ 3\right)  ,\ \left(  2j+3,\ 3\right)  \right\}  $.
This completes the induction step.

Thus, Claim 1 is proved by induction.
\end{proof}

\begin{statement}
\textit{Claim 2:} The number $n$ is even.
\end{statement}

\begin{proof}
[Proof of Claim 2.]Assume the contrary. Thus, $n$ is odd. But $n>1$ (since
$R_{1,3}$ cannot be tiled by dominos), so that $n-1>0$. Hence, $\left(
n-1\right)  /2$ is a positive integer (since $n$ is odd). Therefore, Claim 1
(applied to $i=\left(  n-1\right)  /2$) shows that the tiling $T$ contains the
basement domino $\left\{  \left(  n-2,\ 1\right)  ,\ \left(  n-1,\ 1\right)
\right\}  $, the middle domino $\left\{  \left(  n-1,\ 2\right)  ,\ \left(
n,\ 2\right)  \right\}  $ and the top domino $\left\{  \left(  n-1,\ 3\right)
,\ \left(  n,\ 3\right)  \right\}  $. But $T$ must also include a domino that
contains the square $\left(  n,\ 1\right)  $ (since $\left(  n,\ 1\right)  \in
R_{n,3}$). This domino cannot be vertical (since it would then collide with
the basement domino $\left\{  \left(  n-2,\ 1\right)  ,\ \left(
n-1,\ 1\right)  \right\}  $, which we know to belong to $T$), and cannot be
horizontal either (since it would then collide with the middle domino
$\left\{  \left(  n-1,\ 2\right)  ,\ \left(  n,\ 2\right)  \right\}  $). This
is clearly absurd. This contradiction shows that our assumption was wrong, so
that Claim 2 is proved.

(Alternatively, Claim 2 can be obtained from a parity argument: Since $T$ is a
tiling of $R_{n,3}$ by dominos, the total \# of squares of $R_{n,3}$ must be
even (since each domino covers exactly $2$ squares). But this total \# is
$3n$. Thus, $3n$ must be even, so that $n$ must be even.)
\end{proof}

Now, Claim 2 shows that $n$ is even, so that $n/2$ is a positive integer.
Furthermore, we know that the tiling $T$ contains the left wall, the first
$n-1$ basement dominos
\[
\left\{  \left(  1,1\right)  ,\ \left(  2,1\right)  \right\}
,\ \ \ \ \ \left\{  \left(  3,1\right)  ,\ \left(  4,1\right)  \right\}
,\ \ \ \ \ \ldots,\ \ \ \ \ \left\{  \left(  n-3,\ 1\right)  ,\ \left(
n-2,\ 1\right)  \right\}
\]
(by Claim 1), all the middle dominos
\[
\left\{  \left(  2,2\right)  ,\ \left(  3,2\right)  \right\}
,\ \ \ \ \ \left\{  \left(  4,2\right)  ,\ \left(  5,2\right)  \right\}
,\ \ \ \ \ \ldots,\ \ \ \ \ \left\{  \left(  n-2,\ 2\right)  ,\ \left(
n-1,\ 2\right)  \right\}
\]
(also by Claim 1) and all the top dominos%
\[
\left\{  \left(  2,3\right)  ,\ \left(  3,3\right)  \right\}
,\ \ \ \ \ \left\{  \left(  4,3\right)  ,\ \left(  5,3\right)  \right\}
,\ \ \ \ \ \ldots,\ \ \ \ \ \left\{  \left(  n-2,\ 3\right)  ,\ \left(
n-1,\ 3\right)  \right\}
\]
(also by Claim 1). This leaves only the four squares $\left(  n-1,\ 1\right)
$, $\left(  n,\ 1\right)  $, $\left(  n,\ 2\right)  $ and $\left(
n,\ 3\right)  $ unaccounted for, but there is only one way to tile them:
namely, with the last basement domino $\left\{  \left(  n-1,\ 1\right)
,\ \left(  n,\ 1\right)  \right\}  $ and the right wall $\left\{  \left(
n,2\right)  ,\ \left(  n,3\right)  \right\}  $. Thus, $T$ must contain all the
basement dominos, all the middle dominos, all the top dominos and both walls
(left and right). Since these dominos cover all the squares of $R_{n,3}$, this
entails that $T$ consists precisely of these dominos. In other words,
$T=A_{n}$. The proof of Proposition \ref{prop.gf.weighted-set.domino.Rn3.ABC}
\textbf{(a)} is thus complete. \medskip

\textbf{(b)} Proposition \ref{prop.gf.weighted-set.domino.Rn3.ABC}
\textbf{(b)} is just Proposition \ref{prop.gf.weighted-set.domino.Rn3.ABC}
\textbf{(a)} reflected across the horizontal axis of symmetry of $R_{n,3}$.
\end{proof}
\end{fineprint}
